\section{Ablation and Analysis}
\begin{figure*}[t]
    \centering
    % \hspace{0.2in}MSCOCO~\citep{Chen2015MicrosoftCC} \hspace{2.5in} Flickr30K~\citep{flickrentitiesijcv} \\
    \hspace{0.5in}MSCOCO~\citep{Chen2015MicrosoftCC} \hspace{1in} Flickr30K~\citep{flickrentitiesijcv} \\\includegraphics[width=0.248\textwidth]{fig/coco_umap.pdf}
    \includegraphics[width=0.243\textwidth]{fig/coco_tuned_umap.pdf}  
    \includegraphics[width=0.243\textwidth]{fig/flickr_umap.pdf}
    \includegraphics[width=0.243\textwidth]{fig/flickr_tuned_umap.pdf}
    % LLaVA-OV~\citep{li2024llava}\hspace{1in}\algname \hspace{1in} LLaVA-OV~\citep{li2024llava} \hspace{1in} \algname 
    LLaVA-OV~\citep{li2024llava}\hspace{0.5in}\algname \hspace{0.6in} LLaVA-OV~\citep{li2024llava} \hspace{0.4in} \algname \hspace{0.4in}
    
    \caption{Scatter plots of image-text embeddings from MSCOCO~\citep{Chen2015MicrosoftCC} and Flickr30K~\citep{flickrentitiesijcv} datasets. \algname removes the existing modality gap in multimodal representation space.}
    \label{fig:gap}
    % \vspace{-0.1in}
\end{figure*}
\subsection{Ablation on Training Objectives}
\begin{table}[t]
\centering
\caption{
    Performance comparison across varying weights for language modeling ($\alpha_{\text{lm}}$) and contrastive ($\alpha_{\text{con}}$) objectives.
    %
    Zero-shot retrieval recall@1 is reported for image-to-text (I2T) and text-to-image (T2I) tasks on the Flickr and MSCOCO datasets, alongside multimodal understanding scores on MMMU and MMStar. 
    %
    A ratio of 0 indicates training solely with contrastive loss, while a ratio of 1 implies equal emphasis on both objectives.
}
\setlength{\tabcolsep}{4pt}
% \renewcommand{\arraystretch}{1.1}
% \resizebox{\linewidth}{!}{
\begin{tabular}{ccccccc}
\toprule
\multicolumn{1}{c}{$\alpha_{\text{lm}}:\alpha_{\text{con}}$} 
& \multicolumn{2}{c}{Flickr} 
& \multicolumn{2}{c}{MSCOCO} 
& MMMU & MMStar \\
\cmidrule(lr){2-3} \cmidrule(lr){4-5}
& I2T & T2I & I2T & T2I & & \\
\midrule
0 & 89.0 & 75.0 & 67.0 & 49.6 & 44.4 & 53.9 \\
0.05 & 88.3 & 75.8 & 65.3 & 48.6 & 47.0 & 56.5 \\
0.1 & 87.5 & 75.3 & 64.4 & 47.9 & 47.0 & 56.8 \\
0.2 & 87.0 & 74.7 & 63.3 & 47.0 & 47.4 & 57.3 \\
1 & 84.9 & 73.2 & 61.1 & 45.1 & 47.4 & 57.1 \\
\bottomrule
\end{tabular}
% }
% \vspace*{-\baselineskip}
\label{tab:loss}
\end{table}
%
% \begin{table}[t]
% \centering
% \setlength{\tabcolsep}{5pt}
% \renewcommand{\arraystretch}{1.1}
% \resizebox{\linewidth}{!}{
% \begin{tabular}{ccccccc}
% \toprule
% $\alpha_{\text{lm}}:\alpha_{\text{con}}$ 
% &\multicolumn{2}{c}{Flickr}
% &\multicolumn{2}{c}{MSCOCO} 
% & MMMU & MMStar  \\
% \midrule
% &I2T &T2I &I2T &T2I &Score &Score  \\
% \midrule
% 0 & 89.0 & 75.0 & 67.0 & 49.6 & 44.4 & 53.9 \\
% 0.05 & 88.3 & 75.8 & 65.3 & 48.6 & 47.0 & 56.5 \\
% 0.1 & 87.5 & 75.3 & 64.4 & 47.9 & 47.0 & 56.8 \\
% 0.2 & 87.0 & 74.7 & 63.3 & 47.0 & 47.4 & 57.3 \\
% 1 &	84.9 &	73.2 & 61.1	& 45.1	& 47.4	& 57.1 \\
% \bottomrule
% \end{tabular}
% }
% \caption{
%     Different weights of training objectives. 
%     %
%     Results on zero-shot retrieval recall@1 for both image-to-text (I2T) and text-to-image (T2I) on Flickr and MSCOCO and multimodal understanding on MMMU and MMStar are reported. 
%     %
%     A ratio of 0 means no langage modeling loss, and the model is trained on contrastive loss only. 
% }
% \label{tab:loss}
% \end{table}
% %

We study the effects of two training objectives and compare different variants of the model in \figref{tab:loss}. Given that the scale of the contrastive loss is one-tenth that of the language modeling loss, we experimented with different weightings for the losses, specifically using ratios of  $\alpha_{\text{lm}}:\alpha_{\text{con}}$ = 0, 0.02, 0.1, 0.2, and 1. Among them, $\alpha_{\text{lm}}:\alpha_{\text{con}}$ = 0 is the case where the model is trained using only the contrastive loss and $\alpha_{\text{lm}} = 0$. Results indicate that generally a larger contrastive loss leads to better representation abilities but worse performance on multimodal understanding tasks. We present a trade-off between representation abilities and generation abilities.  

\subsection{Modality Gap}

Modality gap~\citep{liang2022mind} refers to a geometric phenomenon in the representation spaces of multimodal models, where different data modalities, such as images and texts, are embedded at a significant distance from each other, rather than being uniformly distributed as ideally expected. 
This gap, which originates from initialization and persists throughout the contrastive learning process (as seen in models like CLIP~\citep{radford2021learning}), presents a challenge for language-image pretraining by affecting joint data modeling and understanding. Recent studies~\citep{srivastava2024omnivec,oh2024geodesic} suggest that narrowing this gap could improve multimodal representations and enhance performance in downstream tasks. 

We analyze the modality gaps in the multimodal embedding space of the proposed \algname by visualizing the embeddings of 250 randomly selected image-text pairs in MSCOCO and Flickr datasets. Following~\citep{liang2022mind}, we adopt UMAP~\citep{mcinnes2018umap} to reduce the dimensionality of the embeddings to 2 and  
compare the embeddings of the proposed method with LLaVA-OV~\citep{li2024llava} in \figref{fig:gap}.
As shown in the figure, even though LLaVA-OV~\citep{li2024llava} generated both text and image embeddings using the same model, there exists a significant modality gap between text and images. 
In contrast, after finetuning with contrastive-autoregressvie loss, the proposed model, \algname, removes this modality gap in the generated embeddings, where embeddings are aligned based on semantic meaning rather than modality.
%